\subsection{Selector Classification and Fano--Tomotope Convention}
\label{sec:selector-classification}

The \(24\) selector sheets split into \(19\) full-rank sheets and \(5\) defective sheets.
At the mask level the full-rank sheets decompose as
\[
19=12+7.
\]
The robust family is the Type-A co-singleton family
\[
1110,\;1101,\;1011,\;0111,
\]
for which all three residual channels have rank \(81\).  The Type-B adjacent-pair
family contributes only seven full-rank sheets and contains a residual-channel seam.

We therefore fix the project convention by choosing the Type-A mask whose exterior
edge is the tomotope hinge exit.  In the current Fano chart this is
\[
\mathrm{mask}=1110.
\]
The residual channel is then matched to the first diagonal Fano gauge
\[
011\leftrightarrow\mathrm{far},
\]
so the selected sheet is
\[
(1110,r=0).
\]
BT714 verifies that this sheet has rank \(81\), is boundaryless, and maps the chart
\(81\)-sector onto the Levi \(E_4\) Hodge sector.

This is a canonical convention relative to the existing Fano/tomotope chart.  A
stronger uniqueness theorem would still need to quotient all tomotope orientations
and prove that every admissible convention lands in the same \(E_4\) sector.
